% ============================================================
% SECTION 8 — Conclusion
% ============================================================
\section{Conclusion}
\label{sec:conclusion}

This paper introduced GeoRSCT, a diagnostic benchmark that evaluates
flood substrates against the decisions they are used for, not against a
single held-out accuracy score. By organizing evaluation around six
spatial decision geometries---prediction, ranking, clustering, transfer,
relational propagation, and allocation---the benchmark surfaces where a
substrate's evidence is sufficient, where it is provisional, and where
it is absent, before a model is trusted for operational use.

The findings are mixed by design. Prediction is the geometry with the
strongest evidential support: the representation ladder produces
measurable, leakage-controlled skill that improves with added structure.
Ranking and clustering are descriptive: the probes detect reordering
under infrastructure context and residual spatial autocorrelation, but
the verdicts rest on nine cells and carry the uncertainty appropriate to
that sample size. Transfer (leave-event-out retention) is reported with
its gaps. Relational propagation is specified but not yet executed---the
W-matrix ablation variants are designed, and the verdict cannot be issued
until they run. Allocation is necessarily partial: events predating the
temporal sources carry \texttt{not\_applicable\_temporal} annotations,
and the benchmark reports this as an evidence gap, not a defect.

Several limitations bound the findings. The benchmark operates on
$n{=}9$ modelable cells, enough for one pre-registered confirmatory
test but too few for robust cell-level associations (H2b, H4 are
explicitly exploratory). The Riverside scenario, at 172 ZCTA--event
pairs, is the thinnest cell and produces the widest confidence
intervals. The ZIP-as-ZCTA approximation (C-001) is standard practice
but introduces spatial mismatch, particularly in rural areas outside
our study footprints. And the substrate is U.S.-specific: NFIP claims,
FEMA flood zones, and ACS demographics do not transfer to non-U.S.\
contexts without equivalent data infrastructure.

The benchmark's contribution is not a model or a score but a protocol:
a reusable substrate with provenance and field-status annotations, a
six-geometry evaluation framework with pre-committed verdict criteria,
and a worked example---including negative results and evidence
gaps---of what it means to ask whether the data supports the decision,
not just whether the model fits the data.
